\begin{frame}[fragile]{Spread Operator}

\begin{itemize}
	\item We can use an array to aggregate multiple values (lines 1 -- 2).
	\item When we use an array with `...' prepended, we call the `...' a spread operator. 
	\item The spread operator spreads out each element of an array as function arguments (lines 3 -- 4).  
\end{itemize}

\begin{Code}{}%
\begin{lstlisting}[firstnumber=1, label=glabels, xleftmargin=0pt, basicstyle=\scriptsize]% 

var array = ['hello', 'world'];
console.log(array);    // ['hello', 'world']
console.log(...array); // hello world
console.log(array[0], array[1]) // same as ...arrary
\end{lstlisting}%   
\end{Code}%

\end{frame}

\begin{frame}[fragile]{}

\begin{itemize}
	\item Let's say we have a function hello() that requires two arguments. 
	\item When we don't use the spread operator, we should extract the two elements in an array as in line 7.
	\item But we don't have to do it with the spread operator that does the work for us. 
\end{itemize}

\begin{Code}{}%
\begin{lstlisting}[firstnumber=1, label=glabels, xleftmargin=0pt, basicstyle=\scriptsize]% 

var array = ['hello', 'world'];

function hello(x, y) {
  console.log(x);
  console.log(y);
}
hello(array[0], array[1]); // ugly
hello(...array); // better

\end{lstlisting}%   
\end{Code}%

\end{frame}

\begin{frame}[fragile]{}

\begin{itemize}
	\item We can use the spread operator to concatenate arrays into one. 
	\item As the code shows, we can spread two arrays into a new array. 
\end{itemize}

\begin{Code}{}%
\begin{lstlisting}[firstnumber=1, label=glabels, xleftmargin=0pt, basicstyle=\scriptsize]% 

var a = [1,2,3];
var b = [4,5];
var c = [...a, ...b]; // [1,2,3,4,5]
\end{lstlisting}%   
\end{Code}%
\end{frame}

\begin{frame}[fragile]{Spread Operator for an Object}

\begin{itemize}
	\item We can use the spread operator for an object. 
	\item This example shows how we can copy all the elements in an object o1 into o2 with a new element added. 
\end{itemize}

\begin{Code}{}%
\begin{lstlisting}[firstnumber=1, label=glabels, xleftmargin=0pt, basicstyle=\scriptsize]% 

var o1 = { a : 1, b : 2 };
var o2 = { c : 3, ...o1 };
console.log(o2); // { c: 3, a: 1, b: 2 }
\end{lstlisting}%   
\end{Code}%

\end{frame}

\begin{frame}[fragile]{Deep Copy and Shallow Copy}

\begin{itemize}
	\item The spread operator is used for deep copy. 
	\item When we copy arrays, we sometimes make a mistake copying only a reference; we call it a shallow copy (lines 1 - 3). 
	\item To avoid this, we should copy all the elements; we call it a deep copy (lines 5 - 7).
\end{itemize}

\begin{Code}{}%
\begin{lstlisting}[firstnumber=1, label=glabels, xleftmargin=0pt, basicstyle=\scriptsize]% 

var a = [1,2,3];
var b = a; // shallow copy, a and b reference the same array.
b[0] = 1000; console.log(a[0] == b[0]); // true

var a = [1,2,3]
var b = [...a]
b[0] = 2000; console.log(b[0] != a[2])
\end{lstlisting}%   
\end{Code}%
\end{frame}


\begin{frame}[fragile]{Rest Parameter}

\begin{itemize}
	\item The spread operator can be used as a function parameter when we don't know the number of arguments given (lines 1 -- 4).
	\item When we select only some of the arguments, we can use the normal parameter and the spread parameters to select only the arguments into an array (lines 5 -- 7). 
\end{itemize}

\begin{Code}{}%
\begin{lstlisting}[firstnumber=1, label=glabels, xleftmargin=0pt, basicstyle=\tiny]% 

function f(...params){
  console.log(params); 
}
f(1,2,3,4,5,6,7); // [1,2,3,4,5,6,7]
function f2(a,b, ...params){
  console.log(params); 
}
f2(1,2,3,4,5,6,7); // [3,4,5,6,7]
\end{lstlisting}%   
\end{Code}%

\end{frame}

\begin{frame}[fragile]{Sameness Check in JavaScript}

\begin{itemize}
	\item You may remember that in Java, we use `==' operator for primitive type values, and equals() method to compare objects (CSC 260). 
	\item JavaScript uses `==' operator when we compare two values ignoring data type. 
	\item So, as long as values represent the same number, such as 1, 1.0, or even `1'. They are the same with the `==' operator. 
\end{itemize}

\begin{Code}{}%
\begin{lstlisting}[firstnumber=1, label=glabels, xleftmargin=0pt, basicstyle=\scriptsize]% 

console.log(1 == '1') // true
console.log(1 == 1.0) // true
\end{lstlisting}%   
\end{Code}%
\end{frame}

\begin{frame}[fragile]{}

\begin{itemize}
	\item So, when we compare two values and when we need to check the data type also, we should use the `===' or `!==' operator. 
	\item In this example, line 1 returns false.
	\item However, line 2 returns true because there is no integer type in JavaScript; all the numbers are floating-point numbers. 
\end{itemize}

\begin{Code}{}%
\begin{lstlisting}[firstnumber=1, label=glabels, xleftmargin=0pt, basicstyle=\scriptsize]% 

console.log(1 !== '1') // false
console.log(1 !== 1.0) // true
\end{lstlisting}%   
\end{Code}%
\end{frame}

\begin{frame}[fragile]{}

\begin{itemize}
	\item We should also remember that the number 1 and true are the same with the `==' operator but not with the `===' operator. 
	\item Also, 0 and false are the same with `==' operator, but false with `===' operator.
\end{itemize}

\begin{Code}{}%
\begin{lstlisting}[firstnumber=1, label=glabels, xleftmargin=0pt, basicstyle=\scriptsize]% 

console.log(1 == true) // true
console.log(1 === true) // false

console.log(0 == false) // true
console.log(0 === false) // false
\end{lstlisting}%   
\end{Code}%

\end{frame}

\begin{frame}[fragile]{Object comparison}

\begin{itemize}
	\item We discussed the difference between shallow copy (reference copy) and deep copy. 
	\item In this example, the two different objects comparison with `==' and `===' operators return false. 
	\item It is because a and b store each object's reference (address). 
\end{itemize}

\begin{Code}{}%
\begin{lstlisting}[firstnumber=1, label=glabels, xleftmargin=0pt, basicstyle=\scriptsize]% 

var a = { name : 'Sam' };
var b = {...a};

console.log(a == b)  // false
console.log(a === b) // false
\end{lstlisting}%   
\end{Code}%

\end{frame}

\begin{frame}[fragile]{}

\begin{itemize}
	\item To check shallow equality, we can use both `==' and `===' operators.
\end{itemize}

\begin{Code}{}%
\begin{lstlisting}[firstnumber=1, label=glabels, xleftmargin=0pt, basicstyle=\scriptsize]% 

var a = { name : 'Sam' };
var b = a

console.log(a == b)
console.log(a === b)
\end{lstlisting}%   
\end{Code}%
\end{frame}

\begin{frame}[fragile]{}

\begin{itemize}
	\item To check deep equality, we need to compare if the two objects have the same values recursively. 
\end{itemize}

\begin{Code}{}%
\begin{lstlisting}[firstnumber=1, label=glabels, xleftmargin=0pt, basicstyle=\tiny]% 

function deepEqual(object1, object2) {
  const keys1 = Object.keys(object1);
  const keys2 = Object.keys(object2);
  if (keys1.length !== keys2.length) {
    return false;
  }
  for (const key of keys1) {
    const val1 = object1[key];
    const val2 = object2[key];
    const areObjects = isObject(val1) && isObject(val2);
    if (
      areObjects && !deepEqual(val1, val2) ||
      !areObjects && val1 !== val2
    ) {
      return false;
    }
  }
  return true;
}
function isObject(object) {
  return object != null && typeof object === 'object';
}
\end{lstlisting}%   
\end{Code}%

\end{frame}

\begin{frame}[fragile]{null and undefined Comparison}

\begin{itemize}
	\item In JavaScript, null and undefined are different. 
	\item `null' is a value, so it can be defined to some variables (line 1). 
	\item However, `undefined' means a variable is not defined. So, it's indication not a value. 
\end{itemize}

\begin{Code}{}%
\begin{lstlisting}[firstnumber=1, label=glabels, xleftmargin=0pt, basicstyle=\scriptsize]% 

var x = null; 
console.log(x) // null
var y; // y is not defined
console.log(y) // undefined
\end{lstlisting}%   
\end{Code}%

\end{frame}

\begin{frame}[fragile]{}

\begin{itemize}
	\item When we use `==' to compare `null' and `undefined', it returns true. 
	\item We can understand it, as `==' just compares if the two inputs are equivalent; like 1 and '1'.
	\item As we can expect, the `===' returns false when we compare `null' and `undefined.'
\end{itemize}

\begin{Code}{}%
\begin{lstlisting}[firstnumber=1, label=glabels, xleftmargin=0pt, basicstyle=\scriptsize]% 

console.log(null == undefined) // true
console.log(null === undefined) // false
\end{lstlisting}%   
\end{Code}%

\end{frame}
